% ============================================================
%  CHAPTER 10: SIMPLE MACHINES
%  The Physics Codex: A Complete University Foundation
%  Korede Samuel Omotosho (Falcon)
%  Aurikrex Learning Systems, 2026
%
%  File: chapters/part1/ch10_simple_machines.tex
% ============================================================

\chapter{Simple Machines}
\label{ch:simple_machines}
\minitoc
\newpage

\chapterquote{Give me a lever long enough and a fulcrum
on which to place it, and I shall move the world.}{Archimedes}

\begin{objectivebox}
By the end of this chapter, you should be able to:
\begin{enumerate}[noitemsep]
  \item Define mechanical advantage, velocity ratio
  and efficiency of a machine
  \item Explain why efficiency is always less than
  100\% in practice
  \item Analyse the lever and classify lever classes
  with real examples
  \item Derive and apply the velocity ratio for the
  inclined plane, pulley systems, wheel and axle,
  screw jack, gear wheels and wedge
  \item Apply the principle of moments to levers and
  other machines
  \item Solve problems involving overcoming friction
  and calculating actual efficiency
  \item Explain why simple machines are useful even
  when they cannot create energy
\end{enumerate}
\end{objectivebox}

% ============================================================
\section{What is a Machine?}
% ============================================================

A machine is any device that makes work easier to
do. The word ``easier'' here has a precise meaning:
a machine allows us to apply a smaller force over a
larger distance to achieve the same result as a larger
force over a smaller distance. Or it allows us to
change the \textit{direction} of the applied force.

\begin{keybox}[The Fundamental Truth About Machines]
No machine can create energy. A machine can only
\textit{redirect} or \textit{redistribute} the energy
you put into it. In an ideal (frictionless) machine,
energy out equals energy in. In a real machine,
energy out is always \textit{less} than energy in
because some energy is always lost to friction and
heat. \textbf{Machines save effort, not energy.}
\end{keybox}

\begin{historybox}[Archimedes and the Simple Machine]
Archimedes of Syracuse (287--212 BC) was the first
person to give a rigorous mathematical treatment of
simple machines. He analysed the lever, the pulley,
and the screw with extraordinary precision. His
boast about moving the world with a long enough
lever was not mere rhetoric --- it was a mathematically
sound statement about the principle of the lever.
He famously helped defend Syracuse against Roman
ships by designing giant cranes that could lift
ships out of the water using pulley systems, and
catapults that could hurl enormous stones. Engineering
has not fundamentally changed since: we still use
the same simple machines, just combined in more
complex ways.
\end{historybox}

% ============================================================
\section{Key Definitions}
% ============================================================

\begin{definitionbox}[Effort, Load and Mechanical Advantage]
\df{Effort ($E$)}{The force \textit{applied} to the
machine by the operator.}

\df{Load ($L$)}{The force exerted \textit{by} the
machine on the object being moved (the useful output
force).}

\df{Mechanical Advantage (MA)}{The ratio of the load
to the effort:}
\[
  MA = \frac{L}{E}
\]
MA has \textbf{no unit}. $MA > 1$ means the machine
multiplies force (load $>$ effort). This is the most
common reason for using a machine.
\end{definitionbox}

\begin{definitionbox}[Velocity Ratio]
The \textbf{velocity ratio} (VR) is the ratio of the
distance moved by the effort to the distance moved
by the load in the same time:
\[
  VR = \frac{d_E}{d_L}
\]
It is also equal to the ratio of the speed of the
effort point to the speed of the load point.
VR has \textbf{no unit}.

VR is determined purely by the geometry of the
machine --- it does not depend on friction, efficiency,
or the magnitude of the forces.
\end{definitionbox}

\begin{definitionbox}[Efficiency]
The \textbf{efficiency} $\eta$ of a machine is the
ratio of useful work output to total work input:
\[
  \eta = \frac{\text{Work output}}{\text{Work input}}
  \times 100\%
  = \frac{L \times d_L}{E \times d_E} \times 100\%
  = \frac{MA}{VR} \times 100\%
\]
Efficiency is always less than 100\% in real machines
because some energy is always dissipated as heat
through friction. In an ideal (frictionless) machine,
$\eta = 100\%$ and $MA = VR$.
\end{definitionbox}

\begin{keybox}[Relationship Between MA, VR and Efficiency]
\[
  \eta = \frac{MA}{VR} \implies MA = \eta \times VR
\]
For a real machine: $MA < VR$ (because $\eta < 1$)\\
For an ideal machine: $MA = VR$ (because $\eta = 1$)
\end{keybox}

\begin{examplebox}[10.1]
\stepcounter{workedexample}
A machine lifts a $600\ \text{N}$ load through
$2.0\ \text{m}$ when an effort of $150\ \text{N}$ is
applied through $10\ \text{m}$.\\
(a) Calculate the mechanical advantage.\\
(b) Calculate the velocity ratio.\\
(c) Calculate the efficiency.\\
(d) Calculate the work done against friction.
\end{examplebox}

\begin{solutionbox}
\textbf{(a)} Mechanical advantage:
\[
  MA = \frac{L}{E} = \frac{600}{150} = 4.0
\]

\textbf{(b)} Velocity ratio:
\[
  VR = \frac{d_E}{d_L} = \frac{10}{2.0} = 5.0
\]

\textbf{(c)} Efficiency:
\[
  \eta = \frac{MA}{VR} \times 100\%
  = \frac{4.0}{5.0} \times 100\% = 80\%
\]

\textbf{(d)} Work done against friction:
\[
  W_{input} = E \times d_E = 150 \times 10 = 1500\ \text{J}
\]
\[
  W_{output} = L \times d_L = 600 \times 2.0 = 1200\ \text{J}
\]
\[
  W_{friction} = W_{input} - W_{output} = 1500 - 1200
  = 300\ \text{J}
\]
\end{solutionbox}

% ============================================================
\section{The Lever}
% ============================================================

The lever is the simplest machine --- a rigid bar
that pivots about a fixed point called the
\textbf{fulcrum}. The principle behind it is the
\textbf{principle of moments}.

\begin{lawbox}[Principle of Moments]
For a body in equilibrium, the sum of the clockwise
moments about any point equals the sum of the
anticlockwise moments about the same point:
\[
  \sum (\text{clockwise moments})
  = \sum (\text{anticlockwise moments})
\]
For a simple lever with effort arm $d_E$ and load
arm $d_L$:
\[
  E \times d_E = L \times d_L
\]
Therefore:
\[
  MA = \frac{L}{E} = \frac{d_E}{d_L} = VR
  \quad \text{(for an ideal lever)}
\]
\end{lawbox}

\subsection{Classes of Levers}

Levers are classified by the relative positions of
the Effort (E), Fulcrum (F), and Load (L).

\begin{diagbox}[The Three Classes of Lever]
\begin{center}
\begin{tikzpicture}[scale=0.85, font=\small]

  % === FIRST CLASS ===
  \node[font=\bfseries\small, anchor=east] at (-0.2,6.65) {1st};
  % Beam
  \fill[gray!30] (0,6.5) rectangle (7,6.8);
  % Fulcrum (triangle at centre)
  \fill[codexgold] (3.5,5.8) -- (3.0,6.5) -- (4.0,6.5) -- cycle;
  % Effort (left)
  \draw[->, very thick, codexblue] (0.8,6.8) -- (0.8,7.8);
  \node[above, codexblue] at (0.8,7.8) {E};
  % Load (right, down)
  \draw[->, very thick, codexred] (6.2,6.8) -- (6.2,5.8);
  \node[below, codexred] at (6.2,5.8) {L};
  % Labels
  \node[below, gray] at (3.5,5.7) {F};
  \draw[<->, gray] (0.8,5.35) -- (3.5,5.35);
  \node[below, gray] at (2.1,5.25) {$d_E$};
  \draw[<->, gray] (3.5,5.35) -- (6.2,5.35);
  \node[below, gray] at (4.85,5.25) {$d_L$};

  % === SECOND CLASS ===
  \node[font=\bfseries\small, anchor=east] at (-0.2,3.65) {2nd};
  % Beam
  \fill[gray!30] (0,3.5) rectangle (7,3.8);
  % Fulcrum (triangle at left)
  \fill[codexgold] (0.5,2.8) -- (0,3.5) -- (1.0,3.5) -- cycle;
  % Load (centre, down)
  \draw[->, very thick, codexred] (3.5,3.8) -- (3.5,2.8);
  \node[below, codexred] at (3.5,2.8) {L};
  % Effort (right, up)
  \draw[->, very thick, codexblue] (6.5,3.8) -- (6.5,4.8);
  \node[above, codexblue] at (6.5,4.8) {E};
  % Labels
  \node[below, gray] at (0.5,2.7) {F};
  \draw[<->, gray] (0.5,2.35) -- (6.5,2.35);
  \node[below, gray] at (3.5,2.25) {$d_E$};
  \draw[<->, gray] (0.5,1.9) -- (3.5,1.9);
  \node[below, gray] at (2.0,1.8) {$d_L$};

  % === THIRD CLASS ===
  \node[font=\bfseries\small, anchor=east] at (-0.2,0.65) {3rd};
  % Beam
  \fill[gray!30] (0,0.5) rectangle (7,0.8);
  % Fulcrum (triangle at left)
  \fill[codexgold] (0.5,-0.2) -- (0,0.5) -- (1.0,0.5) -- cycle;
  % Effort (centre, up)
  \draw[->, very thick, codexblue] (2.5,0.8) -- (2.5,1.8);
  \node[above, codexblue] at (2.5,1.8) {E};
  % Load (right, down)
  \draw[->, very thick, codexred] (6.2,0.8) -- (6.2,-0.2);
  \node[below, codexred] at (6.2,-0.2) {L};
  % Labels
  \node[below, gray] at (0.5,-0.3) {F};
  \draw[<->, gray] (0.5,-0.8) -- (2.5,-0.8);
  \node[below, gray] at (1.5,-0.9) {$d_E$};
  \draw[<->, gray] (0.5,-1.3) -- (6.2,-1.3);
  \node[below, gray] at (3.35,-1.4) {$d_L$};
  
\end{tikzpicture}
\end{center}
\end{diagbox}

\begin{defbox}[Lever Classes --- Summary]
\begin{center}
\begin{tabular}{lllll}
\toprule
\textbf{Class} & \textbf{Arrangement} & \textbf{MA} &
\textbf{Examples}\\
\midrule
1st & F between E and L & Can be  $\gtreqqless 1$ & See-saw, crowbar, scissors, pliers\\
2nd & L between F and E & Always $>1$ &
Wheelbarrow, nutcracker, bottle opener\\
3rd & E between F and L & Always $<1$ &
Tweezers, forearm, fishing rod, broom\\
\bottomrule
\end{tabular}
\end{center}
\end{defbox}

\begin{realworldbox}[The Human Forearm as a Third-Class Lever]
Your forearm is a third-class lever. The fulcrum is
the elbow joint, the effort is the biceps muscle
attached close to the elbow, and the load is whatever
you are holding in your hand. Since the effort arm
is much shorter than the load arm, $MA < 1$ --- the
muscle must exert a much larger force than the weight
of the object being held. When you hold a $1\ \text{kg}$
book in your hand ($\approx 30\ \text{cm}$ from elbow)
and your biceps attaches $4\ \text{cm}$ from the elbow,
the biceps must exert $30/4 = 7.5$ times the weight
of the book --- about $75\ \text{N}$ to hold a $10\ \text{N}$
load. The advantage of third-class levers is speed
and range of movement, not force multiplication.
\end{realworldbox}

\begin{examplebox}[10.2]
\stepcounter{workedexample}
A crowbar (first-class lever) has its fulcrum
$15\ \text{cm}$ from the load and $90\ \text{cm}$
from where the effort is applied. A $720\ \text{N}$
load is to be lifted.\\
(a) Calculate the velocity ratio.\\
(b) Calculate the effort required if the crowbar
is 90\% efficient.\\
(c) Calculate the mechanical advantage.
\end{examplebox}

\begin{solutionbox}
\textbf{(a)} VR (for a lever):
\[
  VR = \frac{d_E}{d_L} = \frac{90}{15} = 6.0
\]

\textbf{(b)} From $\eta = MA/VR$:
\[
  MA = \eta \times VR = 0.90 \times 6.0 = 5.4
\]
From $MA = L/E$:
\[
  E = \frac{L}{MA} = \frac{720}{5.4}
  = 133.3\ \text{N} \approx 133\ \text{N}
\]

\textbf{(c)} Mechanical advantage already found:
$MA = 5.4$
\end{solutionbox}

% ============================================================
\section{The Inclined Plane}
% ============================================================

\begin{processbox}[Inclined Plane Analysis]
An inclined plane allows a heavy object to be raised
to a height $h$ by pushing it along a longer slope
of length $\ell$, reducing the required force.

\begin{center}
\begin{tikzpicture}[scale=1.0]
  % Incline
  \fill[gray!20] (0,0) -- (6,0) -- (6,3) -- cycle;
  \draw[thick] (0,0) -- (6,0) -- (6,3) -- (0,0);

  % Object
  \begin{scope}[rotate around={26.6:(2.5,1.25)}]
    \fill[codexblue!40] (2.0,1.1) rectangle (3.0,1.6);
  \end{scope}
  \node[font=\small] at (2.5,1.4) {$m$};

  % Effort arrow (along slope)
  \draw[->, very thick, codexred]
    (2.5,1.4) -- ++(1.5,0.75)
    node[above, font=\small]{$E$};

  % Height and length labels
  \draw[<->, gray] (6.3,0) -- (6.3,3);
  \node[right, gray] at (6.4,1.5) {$h$};
  \draw[<->, gray] (0,-0.4) -- (6,-0.4);
  \node[below, gray] at (3,-0.4) {$\ell \sin\theta = h$};

  % Slope length
  \draw[<->, codexgold] (-0.2,0) -- (5.8,2.9);
  \node[left, codexgold] at (2.5,1.8) {$\ell$};

  % Angle
  \draw[codexgold] (0.8,0) arc(0:26.6:0.8);
  \node[codexgold] at (1.2,0.18) {$\theta$};
\end{tikzpicture}
\end{center}

The effort moves through distance $\ell$ (length of
slope) while the load is raised by $h$ (vertical height):
\[
  VR = \frac{d_E}{d_L} = \frac{\ell}{h}
  = \frac{1}{\sin\theta}
\]
For an ideal inclined plane (no friction):
\[
  E \times \ell = L \times h
  \implies E = L\sin\theta
\]
The shallower the incline, the smaller the effort
required, but the longer the distance pushed.
\end{processbox}

\begin{examplebox}[10.3]
\stepcounter{workedexample}
A ramp is $5.0\ \text{m}$ long and raises a load
to a height of $1.5\ \text{m}$.
A $200\ \text{N}$ effort is used to push a
$600\ \text{N}$ load up the ramp.\\
(a) Find the VR and the angle of the incline.\\
(b) Find the efficiency.\\
(c) Find the friction force acting on the load.
\end{examplebox}

\begin{solutionbox}
\textbf{(a)} VR and angle:
\[
  VR = \frac{\ell}{h} = \frac{5.0}{1.5} = 3.33
\]
\[
  \sin\theta = \frac{h}{\ell} = \frac{1.5}{5.0} = 0.30
  \implies \theta = 17.5°
\]

\textbf{(b)} MA $= L/E = 600/200 = 3.0$
\[
  \eta = \frac{MA}{VR} \times 100\%
  = \frac{3.0}{3.33} \times 100\% = 90.1\%
\]

\textbf{(c)} Work input $= 200 \times 5.0 = 1000\ \text{J}$\\
Work output $= 600 \times 1.5 = 900\ \text{J}$\\
Work lost to friction $= 1000 - 900 = 100\ \text{J}$\\
\[
  f = \frac{W_{friction}}{d_E}
  = \frac{100}{5.0} = 20\ \text{N}
\]
\end{solutionbox}

% ============================================================
\section{Pulleys}
% ============================================================

A pulley is a wheel with a grooved rim over which a
rope or chain runs. A single fixed pulley only changes
the \textit{direction} of force, not its magnitude
($VR = 1$). Combining multiple pulleys reduces the
required effort.

\begin{lawbox}[Pulley Systems]
For a pulley system, the VR equals the number of
rope segments \textbf{supporting the load} (i.e.\
the number of sections of rope attached to or
supporting the lower (movable) block):
\[
  VR = n \quad \text{(number of load-supporting ropes)}
\]
For an ideal pulley system:
\[
  MA = VR = n
  \implies E = \frac{L}{n}
\]
\end{lawbox}

\begin{diagbox}[Pulley Systems]
\begin{center}
\begin{tikzpicture}[scale=0.8, font=\small]

  % SINGLE FIXED PULLEY (VR = 1)
  \node[font=\bfseries, align=center] at (1.5,10.0) {Single fixed\\($VR=1$)};
  \draw[thick, fill=gray!30] (1.5,7.5) circle (0.5cm);
  \fill (1.5,7.5) circle (2pt);
  % Rope
  \draw[thick] (1.0,7.5) -- (1.0,6.0);
  \draw[thick] (2.0,7.5) -- (2.0,9.0);
  % Load
  \fill[codexblue!40] (0.6,5.5) rectangle (1.4,6.0);
  \node at (1.0,5.75) {$L$};
  % Effort
  \draw[->, thick, codexred] (2.0,8.5) -- (2.0,9.0)
    node[right]{$E=L$};
  % Ceiling
  \fill[gray!30] (1.0,7.8) rectangle (2.0,8.1);
  \draw[thick] (1.0,7.8) -- (2.0,7.8);

  % MOVABLE PULLEY SYSTEM (VR = 2)
  \node[font=\bfseries, align=center] at (6,10.0) {One movable\\($VR=2$)};
  % Fixed pulley
  \draw[thick, fill=gray!30] (6,7.5) circle (0.4cm);
  \fill (6,7.5) circle (2pt);
  % Movable pulley
  \draw[thick, fill=gray!30] (6,5.8) circle (0.4cm);
  \fill (6,5.8) circle (2pt);
  % Ropes
  \draw[thick] (5.6,5.8) -- (5.6,7.5);  % left rope
  \draw[thick] (6.4,5.8) -- (6.4,7.0);  % right rope (attached to ceiling)
  \draw[thick] (5.6,7.5) -- (6.0,7.1);  % over fixed pulley
  \draw[thick] (6.0,7.9) -- (6.0,8.6);  % effort rope
  % Ceiling
  \fill[gray!30] (5.0,7.8) rectangle (7.0,8.1);
  \draw[thick] (5.0,7.8) -- (7.0,7.8);
  \draw[thick] (6.4,7.8) -- (6.4,7.0); % right support to ceiling
  % Load
  \fill[codexblue!40] (5.4,5.2) rectangle (6.6,5.6);
  \node at (6.0,5.4) {$L$};
  % Effort
  \draw[->, thick, codexred]
    (6.0,8.3) -- (6.0,9.0) node[right]{$E = L/2$};
  % Labels
  \node[codexblue, font=\tiny] at (5.2,6.6) {1};
  \node[codexblue, font=\tiny] at (6.65,6.6) {2};

  % VR labels
  \node[codexgreen, font=\small] at (1.5,5.0)
    {$VR = 1$};
  \node[codexgreen, font=\small] at (6.0,4.8)
    {$VR = 2$};
\end{tikzpicture}
\end{center}
\end{diagbox}

\begin{examplebox}[10.4]
\stepcounter{workedexample}
A pulley system has 4 load-supporting ropes.
An effort of $250\ \text{N}$ lifts a load of
$800\ \text{N}$ through $2.0\ \text{m}$.\\
(a) Find the VR.\\
(b) Find the MA.\\
(c) Find the efficiency.\\
(d) Find the effort rope distance moved.
\end{examplebox}

\begin{solutionbox}
\textbf{(a)} $VR = 4$ (four load-supporting ropes)

\textbf{(b)} $MA = L/E = 800/250 = 3.2$

\textbf{(c)} $\eta = MA/VR \times 100\%
= 3.2/4 \times 100\% = 80\%$

\textbf{(d)} From $VR = d_E/d_L$:
\[
  d_E = VR \times d_L = 4 \times 2.0 = 8.0\ \text{m}
\]
\end{solutionbox}

% ============================================================
\section{The Wheel and Axle}
% ============================================================

\begin{processbox}[Wheel and Axle]
A wheel of radius $R$ is mounted on an axle of
radius $r$. The effort is applied at the rim of
the wheel, and the load is lifted by a rope wound
around the axle.

The effort moves through a circumference of the
wheel, the load moves through a circumference of
the axle (in one turn):
\[
  VR = \frac{2\pi R}{2\pi r} = \frac{R}{r}
\]
Examples: steering wheel and steering column;
winch handle and drum; screwdriver handle and shaft.
\end{processbox}

\begin{examplebox}[10.5]
\stepcounter{workedexample}
A well windlass (wheel and axle) has a wheel
radius of $35\ \text{cm}$ and an axle radius of
$5\ \text{cm}$. An effort of $80\ \text{N}$ raises
a $400\ \text{N}$ bucket of water.\\
(a) Find the VR.\\
(b) Find the MA.\\
(c) Find the efficiency.
\end{examplebox}

\begin{solutionbox}
\textbf{(a)} $VR = R/r = 35/5 = 7.0$

\textbf{(b)} $MA = L/E = 400/80 = 5.0$

\textbf{(c)} $\eta = MA/VR \times 100\%
= 5.0/7.0 \times 100\% = 71.4\%$

\begin{realworldbox}[The Well Windlass in Rural Nigeria]
The traditional hand-operated well windlass is still
widely used for drawing water in rural Nigerian
communities, particularly in the North where deep
wells are common. The windlass is a direct application
of the wheel-and-axle principle: the handle is the
``wheel'' (large radius), and the drum around which
the rope winds is the ``axle'' (small radius).
The mechanical advantage allows a woman or child to
lift a $20\ \text{kg}$ bucket of water with a fraction
of the force that would be required without the machine.
\end{realworldbox}
\end{solutionbox}

% ============================================================
\section{The Screw Jack}
% ============================================================

\begin{processbox}[Screw Jack]
A screw jack converts rotational motion into linear
motion. The effort is applied at the end of a handle
of length $r$, and the load advances by one
\textbf{pitch} $p$ (the distance between adjacent
threads) per complete rotation.

Per one complete turn of the handle:
\begin{itemize}[noitemsep]
  \item Effort moves through: $2\pi r$ (circumference
  of circle traced by handle)
  \item Load moves through: $p$ (one pitch)
\end{itemize}
\[
  VR = \frac{2\pi r}{p}
\]
\end{processbox}

\begin{examplebox}[10.6]
\stepcounter{workedexample}
A screw jack has a handle of length $30\ \text{cm}$
and a pitch of $5.0\ \text{mm}$. An effort of
$50\ \text{N}$ is applied.\\
(a) Find the VR.\\
(b) Find the load that can be lifted if the
efficiency is $35\%$.
\end{examplebox}

\begin{solutionbox}
\textbf{(a)}
\[
  VR = \frac{2\pi r}{p}
  = \frac{2\pi \times 0.30}{0.005}
  = \frac{1.885}{0.005} = 377
\]

\textbf{(b)}
\[
  MA = \eta \times VR = 0.35 \times 377 = 132
\]
\[
  L = MA \times E = 132 \times 50 = 6600\ \text{N}
\]

A $50\ \text{N}$ effort lifts $6600\ \text{N}$
($\approx 673\ \text{kg}$). This is why car jacks
can lift a vehicle: the screw mechanism provides
an enormous VR.
\end{solutionbox}

% ============================================================
\section{Gear Wheels}
% ============================================================

\begin{processbox}[Gear Wheels]
When two meshed gear wheels (driver and driven)
rotate together, each tooth on the driver engages
one tooth on the driven wheel. The number of teeth
on each wheel determines the speed ratio.

If the driver gear has $N_1$ teeth and the driven
gear has $N_2$ teeth:
\[
  VR = \frac{N_2}{N_1}
  = \frac{\text{teeth on driven gear}}
  {\text{teeth on driver gear}}
\]
If $N_2 > N_1$: driven gear rotates more slowly but
with more torque (force amplification).\\
If $N_2 < N_1$: driven gear rotates faster but with
less torque (speed amplification).
\end{processbox}

\begin{realworldbox}[Gears in Cars and Okada Bikes]
Every bicycle and motorcycle (okada) on Nigerian
roads uses gear systems. In low gear (large VR),
the wheel turns slowly but with great force ---
ideal for climbing hills. In high gear (small VR),
the wheel turns fast with less force --- ideal for
flat roads at speed. The rider selects the gear
that matches the terrain, optimising the trade-off
between speed and force --- exactly the principle
of simple machines.
\end{realworldbox}

% ============================================================
\section{The Wedge}
% ============================================================

\begin{processbox}[Wedge]
A wedge is a double inclined plane. The effort
is applied to the blunt end and the wedge is
driven forward, splitting or lifting the load
sideways.

For a wedge of length $\ell$ and thickness $t$:
\[
  VR = \frac{\ell}{t}
  = \frac{\text{slant length}}{\text{thickness}}
\]
A thin, sharp wedge has a high VR and requires
less effort but must be driven further.
Applications: axes, chisels, knives, ploughs,
door wedges.
\end{processbox}

% ============================================================
\section{Summary Table of Simple Machines}
% ============================================================

\begin{table}[H]
\centering
\caption{Summary of Simple Machines and Their Velocity Ratios}
\label{tab:machines_summary}
\begin{tabular}{lll}
\toprule
\textbf{Machine} & \textbf{VR Formula} & \textbf{Example}\\
\midrule
Lever (1st class) &
  $\dfrac{d_E}{d_L}$ &
  Crowbar, scissors\\[8pt]
Lever (2nd class) &
  $\dfrac{d_E}{d_L} > 1$ (always) &
  Wheelbarrow\\[8pt]
Lever (3rd class) &
  $\dfrac{d_E}{d_L} < 1$ (always) &
  Tweezers, forearm\\[8pt]
Inclined plane &
  $\dfrac{\ell}{h} = \dfrac{1}{\sin\theta}$ &
  Ramp, road cutting\\[8pt]
Pulley system &
  $n$ (number of support ropes) &
  Block and tackle\\[8pt]
Wheel and axle &
  $\dfrac{R}{r}$ &
  Windlass, steering wheel\\[8pt]
Screw jack &
  $\dfrac{2\pi r}{p}$ &
  Car jack, bolt\\[8pt]
Gear wheels &
  $\dfrac{N_2}{N_1}$ &
  Bicycle, gearbox\\[8pt]
Wedge &
  $\dfrac{\ell}{t}$ &
  Axe, chisel, knife\\
\bottomrule
\end{tabular}
\end{table}

% ============================================================
\section{Chapter Summary}
% ============================================================

\begin{summarybox}
\textbf{Mechanical Advantage:} $MA = L/E$ (no unit)

\textbf{Velocity Ratio:} $VR = d_E/d_L$ (no unit,
geometry only)

\textbf{Efficiency:}
$\eta = (MA/VR) \times 100\% < 100\%$ (always)

\textbf{No machine creates energy:}
$W_{out} \leq W_{in}$ always.

\textbf{Lever classes:}
1st (FEL between E-F-L);
2nd (L between E and F, $MA > 1$ always);
3rd (E between F and L, $MA < 1$ always)

\textbf{Principle of moments:}
$E \times d_E = L \times d_L$ (ideal lever)

\textbf{VR formulae:}
Inclined plane: $1/\sin\theta$;
Pulley: $n$;
Wheel and axle: $R/r$;
Screw jack: $2\pi r/p$;
Gears: $N_2/N_1$;
Wedge: $\ell/t$
\end{summarybox}

% ============================================================
\newpage
\begin{exercisebox}[10]

\subsection*{Section A: Multiple Choice}

\textbf{A1.} A machine has $VR = 6$ and efficiency
$= 75\%$. Its mechanical advantage is:

\mcoption{A}{$8.0$}
\mcoption{B}{$4.5$}
\mcoption{C}{$6.0$}
\mcoption{D}{$0.75$}
\ansref{Ch10-A1}

\vspace{0.4cm}

\textbf{A2.} Which of the following always has
$MA < 1$?

\mcoption{A}{First class lever}
\mcoption{B}{Second class lever}
\mcoption{C}{Third class lever}
\mcoption{D}{Inclined plane}
\ansref{Ch10-A2}

\vspace{0.4cm}

\textbf{A3.} A pulley system uses 5 load-supporting
rope segments. An ideal machine of this type lifts
a $1000\ \text{N}$ load with an effort of:

\mcoption{A}{$50\ \text{N}$}
\mcoption{B}{$100\ \text{N}$}
\mcoption{C}{$200\ \text{N}$}
\mcoption{D}{$500\ \text{N}$}
\ansref{Ch10-A3}

\vspace{0.4cm}

\textbf{A4.} A screw jack has pitch $4\ \text{mm}$
and handle length $20\ \text{cm}$.
The velocity ratio is (take $\pi = 3.14$):

\mcoption{A}{$50$}
\mcoption{B}{$157$}
\mcoption{C}{$314$}
\mcoption{D}{$628$}
\ansref{Ch10-A4}

\vspace{0.4cm}

\textbf{A5.} An inclined plane is $8\ \text{m}$ long
and $2\ \text{m}$ high. Its velocity ratio is:

\mcoption{A}{$2$}
\mcoption{B}{$4$}
\mcoption{C}{$6$}
\mcoption{D}{$16$}
\ansref{Ch10-A5}

\vspace{0.4cm}

\textbf{A6.} In a wheel-and-axle system, the wheel
radius is $40\ \text{cm}$ and the axle radius is
$8\ \text{cm}$. To lift a $600\ \text{N}$ load with
an ideal machine, the effort required is:

\mcoption{A}{$75\ \text{N}$}
\mcoption{B}{$120\ \text{N}$}
\mcoption{C}{$150\ \text{N}$}
\mcoption{D}{$240\ \text{N}$}
\ansref{Ch10-A6}

\vspace{0.4cm}

\textbf{A7.} Which statement about machines is
correct?

\mcoption{A}{A machine can create energy if designed
carefully}
\mcoption{B}{The efficiency of an ideal machine is
always $100\%$}
\mcoption{C}{$MA$ can exceed $VR$ in a real machine}
\mcoption{D}{A single fixed pulley has $MA > 1$}
\ansref{Ch10-A7}

\vspace{0.4cm}

\textbf{A8.} A wheelbarrow (second class lever) has
its wheel (fulcrum) $0.3\ \text{m}$ from the load
and handles $1.5\ \text{m}$ from the load. The MA is:

\mcoption{A}{$0.2$}
\mcoption{B}{$1.2$}
\mcoption{C}{$5.0$}
\mcoption{D}{$6.0$}
\ansref{Ch10-A8}

\vspace{0.4cm}

\textbf{A9.} A gear system has a driver gear with
$20$ teeth and a driven gear with $80$ teeth.
If the driver rotates at $400\ \text{rpm}$, the
driven gear rotates at:

\mcoption{A}{$1600\ \text{rpm}$}
\mcoption{B}{$400\ \text{rpm}$}
\mcoption{C}{$100\ \text{rpm}$}
\mcoption{D}{$80\ \text{rpm}$}
\ansref{Ch10-A9}

\vspace{0.4cm}

\textbf{A10.} The efficiency of a machine is $60\%$
and its $VR = 5$. An effort of $200\ \text{N}$ is
applied. The load that can be lifted is:

\mcoption{A}{$240\ \text{N}$}
\mcoption{B}{$400\ \text{N}$}
\mcoption{C}{$600\ \text{N}$}
\mcoption{D}{$1000\ \text{N}$}
\ansref{Ch10-A10}

\vspace{0.6cm}

\subsection*{Section B: Theory and Calculations}

\textbf{B1.} A first-class lever has its fulcrum
$20\ \text{cm}$ from a $500\ \text{N}$ load and
$80\ \text{cm}$ from where the effort is applied.
($g = 10\ \text{m\,s}^{-2}$)\\[4pt]
(a) Calculate the ideal (frictionless) effort needed.\\
(b) The actual effort required is $140\ \text{N}$.
Calculate the actual MA and efficiency.\\
(c) Calculate the work done against friction when the
load is raised by $0.5\ \text{m}$.\\
(d) A person of mass $60\ \text{kg}$ stands on one
end of a $3.0\ \text{m}$ see-saw (first-class lever)
with the fulcrum at the centre. At what distance from
the fulcrum must an $80\ \text{kg}$ person sit on the
other side for the see-saw to balance?
\hfill\ansref{Ch10-B1}

\vspace{0.4cm}

\textbf{B2.} A block-and-tackle pulley system has
3 pulleys in the upper (fixed) block and 2 pulleys
in the lower (movable) block. A $1800\ \text{N}$
load is lifted $3.0\ \text{m}$ by an effort of
$420\ \text{N}$.\\[4pt]
(a) State the number of load-supporting rope segments
and hence the VR.\\
(b) Calculate the MA and efficiency.\\
(c) Calculate the effort distance moved.\\
(d) Calculate the energy wasted due to friction during
this lift.\\
(e) If the efficiency could be improved to $90\%$,
what effort would be needed to lift the same load?
\hfill\ansref{Ch10-B2}

\vspace{0.4cm}

\textbf{B3.} A screw jack is used to lift a car of
mass $1200\ \text{kg}$ off the ground.
The jack has a handle of length $25\ \text{cm}$
and a pitch of $4\ \text{mm}$.
($g = 10\ \text{m\,s}^{-2}$, $\pi = 3.14$)\\[4pt]
(a) Calculate the weight of the car.\\
(b) Calculate the VR of the screw jack.\\
(c) If the efficiency is $30\%$, calculate the MA.\\
(d) Calculate the effort required to lift the car.\\
(e) How far must the handle move to raise the car
by $10\ \text{cm}$? Express in metres.\\
(f) Explain why the efficiency of a screw jack is
so low and why this low efficiency is actually
useful in practice.
\hfill\ansref{Ch10-B3}

\vspace{0.4cm}

\textbf{B4.} A bicycle has a pedal gear with
$40$ teeth and a rear wheel sprocket with $16$ teeth.
The rear wheel has a diameter of $70\ \text{cm}$.
The cyclist applies a force of $120\ \text{N}$ on
each pedal. The pedal crank radius is $18\ \text{cm}$.\\[4pt]
(a) Calculate the velocity ratio of the gear system.\\
(b) How far does the bicycle travel per complete
revolution of the pedals?\\
(c) If the efficiency of the chain-gear system is
$95\%$, calculate the driving force at the rear
wheel rim.\\
(d) At a pedalling rate of $80\ \text{rpm}$,
calculate the cyclist's speed in km\,h$^{-1}$.\\
(e) Calculate the power output of the cyclist.
\hfill\ansref{Ch10-B4}

\vspace{0.4cm}

\textbf{B5.} A student investigates a pulley system
and records the following data:

\begin{center}
\begin{tabular}{ccc}
\toprule
Load (N) & Effort (N) & Efficiency (\%)\\
\midrule
100 & 35 & 57.1\\
200 & 60 & 66.7\\
300 & 85 & 70.6\\
400 & 110 & 72.7\\
500 & 135 & 74.1\\
\bottomrule
\end{tabular}
\end{center}

The system has $VR = 5$.\\[4pt]
(a) Verify the VR from the first row of data and
identify why the efficiency is not $100\%$.\\
(b) Explain the trend: why does efficiency increase
as the load increases?\\
(c) Plot effort against load and use the graph
to find:\\
\quad (i) The effort needed to overcome friction
alone (y-intercept).\\
\quad (ii) The effort needed to lift a $700\ \text{N}$
load (extrapolation).\\
(d) Derive the relationship:
$E = \dfrac{L}{\eta \times VR}$ and show that for
this system $E = aL + b$ (linear), identifying
the physical meaning of $a$ and $b$.
\hfill\ansref{Ch10-B5}

\end{exercisebox}
